%! TeX program = xelatex
% Resume in LaTeX. Author: Martin Yang. License: MIT.
\documentclass[a4paper,10pt]{article}
\usepackage{xeCJK}
\usepackage{resume-layout}
\usepackage{fontawesome}
\renewcommand{\resumeHeaderIcon}[1]{#1\hspace{3pt}}
\setmainfont{Noto Serif}
\setCJKmainfont{Noto Serif CJK SC}
\titlespacing*{\section}{0pt}{8pt}{4pt}
\hypersetup{pdfauthor={Yang Ruoqi},pdftitle={杨若琦 - 中文简历},pdflang={zh-CN}}

\begin{document}

\resumeHeader{杨若琦}{广东\enspace 深圳}
  {github.com/yangrq1018}{yangruoqi-site.pages.dev}

\section{教育经历}

\resumeEntry{伦敦帝国理工}{英国\enspace 伦敦}
  {金融学硕士\quad\textbf{毕业成绩: Distinction}}{2019.09 -- 2020.06}

\resumeEntry{香港大学}{中国\enspace 香港}
  {经济\&金融学学士; 辅修计算机科学; \textbf{GPA: 3.65/4.30}}{2015.09 -- 2019.06}

\resumeEntry{乔治华盛顿大学}{美国\enspace 华盛顿特区}
  {交换生}{2018.01 -- 2018.06}

\section{工作经历}

\resumeEntry{银泰证券证投部}{中国\enspace 深圳}
  {投资经理}{2026.01 -- 2026.11}
\begin{resumeItems}
  \item 负责银泰003转债1号产品投资管理；2026.05.19--10.30期间，相对转债ETF（511380.SH）累计\textbf{超额6.51个百分点}，周度净值最大回撤\textbf{7.92\%}（ETF\textbf{14.43\%}）。
  \item 负责自营产品资金配置，分配组内股票、转债、股指期货策略及货币基金类资产的权重
\end{resumeItems}

\resumeEntry{小牛私募基金管理有限公司}{中国\enspace 深圳}
  {基金经理}{2024.06 -- 2026.01}
\begin{resumeItems}
  \item 实盘进一步优化过的可转债多因子策略，负责对接XTP交易系统的开发和上交所机房部署工作，进行策略的Barra模型风险分析
  \item 搭建一套端到端投研体系，涵盖数据获取、入库、因子研究、机器学习模型、回测、实盘交易以及\mbox{投后评估等模块}
  \item 使用逐笔、委托数据，丰富股票多因子策略的因子库
\end{resumeItems}

\resumeEntry{九坤投资}{中国\enspace 上海}
  {量化研究员}{2022.08 -- 2023.06}
\begin{resumeItems}
  \item 负责可转债策略的实盘维护工作，继续可转债领域的研究
  \item 在多因子框架下，负责股票信号研究，频率包括日频、日内
  \item 参与组内使用的各种多因子研究工具的开发和维护
\end{resumeItems}

\resumeEntry{上海悬铃私募基金管理有限公司}{中国\enspace 上海}
  {量化研究员}{2020.06 -- 2022.07}
\begin{resumeItems}
  \item 研究并落地可转债定价中性策略，以IC股指期货对冲，实盘运行7个月，\textbf{最高管理规模九千万元}
  \item 使用Golang开发高效交易执行程序，向交易网关挂单
  \item 负责QR研究数据库开发，从数据供应商API获取、清理并持久化数据
\end{resumeItems}

\section{奖项与成就}

\resumeFact{普通高等学校招生统一考试：}{广东省前200名}
\resumeFact{CFA Level 1考试：}{通过}
\resumeFact{李兆基奖学金：}{颁发给大三学生中学业成绩优秀者}
\resumeFact{两次院长荣誉榜：}{在第二和第三学年获得}
\resumeFact{国际英语语言测试系统（IELTS）：}{总分8分}
\resumeFact{托福考试（TOEFL）：}{总分106分}

\section{语言和技能}

\resumeFact{计算机：}{掌握C++, Python, Golang; 熟练使用Linux系统}
\resumeFact{语言：}{普通话（母语）, 英语（流利）, 广东话（日常交流）}
\resumeFact{开源贡献：}{向多个超过1000星的Github知名项目贡献过代码，热衷于开源事业}

\end{document}
